\section{Conclusion}\label{sec:conclusion}

We presented ColliForce-VLA, a method that integrates differentiable workspace Signed Distance Fields into the forward pass of a Vision-Language-Action model for collision-aware action generation. By attaching a lightweight SDF Decoder and End-Effector Extractor to the action expert, and training LoRA parameters to reproduce SDF-gradient-corrected action targets, the method internalizes collision avoidance into the model's own parameters. The default inference path requires only the base model with trained LoRA, adding no extra modules or computation. On the SafeLIBERO benchmark, ColliForce-VLA improves Collision Avoidance Rate by \textbf{[TODO]} percentage points over the LoRA baseline while maintaining competitive Task Success Rates, and ablation studies confirm that the gradient-guided refinement contributes beyond auxiliary SDF supervision alone.

\paragraph{Limitations.} Our evaluation is limited to two SafeLIBERO scenes in simulation; broader task coverage and real-robot transfer remain to be demonstrated. The SDF ground truth requires known obstacle meshes, which are available in simulation but would require reconstruction methods (e.g., neural SDF from depth observations) for real-world deployment. The method is validated on a single base VLA ($\pi_0$); generalization to other architectures (OpenVLA, RT-2) is an open question.

\paragraph{Future work.} Three directions follow naturally. First, replacing precomputed SDF with a real-time neural SDF estimated from depth observations would remove the dependency on known meshes and enable deployment in unstructured environments. Second, real-robot validation would test whether the collision avoidance behavior learned in simulation transfers across the sim-to-real gap. Third, extending the evaluation to a broader set of SafeLIBERO scenarios and additional benchmarks would strengthen the generality claims.